\section{Evaluation}

\label{sec:evaluation}

\issue{this need rewriting}
In this section, we first conduct a set of microbenchmarks on the reference monitor and \libcapac to illustrate the overhead induced by individual operations of \thename (\cref{sec:evaluation:mon-microbench}).
% We measure the performance overheads of \thename's instrumentation on the nbench-byte benchmark to study its impact on real-world applications~(\cref{sec:evaluation:nbench}).
We describe the adaptation of \thename in protecting sensitive resources of three real-world applications in \cref{sec:evaluation:porting-app}.
Finally, we evaluate the performance of \thename-adapted applications in \cref{sec:evaluation:macrobench} to show \thename's overall impact on performance.
% of throughput of \thename-enabled webserver, one of the adapted applications,  

% Nevertheless, we plan to evaluate \thename on hardware with native support for \gls{mte} in the future.
\hdr{Evaluation method} We ensured the functional correctness of our implementation by using a QEMU~\cite{qemu} ARM virtual machine that supports both \gls{pa} and \gls{mte} and conducted performance evaluations using the Apple Mac Mini with an M1 processor running Asahi Linux~\cite{asahi-linux}.
The M1 processor includes the \gls{pa} extension; however, to our knowledge, there is no publicly available hardware with ARM MTE.
The QEMU's MTE emulation is mature to a point where it is used for developing a \gls{mte}-based security feature for the Linux kernel to be used in the near future~\cite{kasan}.
We follow the method from a previous work~\cite{hakc} that emulated the \emph{worst-case} performance impact of \gls{mte} by using \emph{\gls{mte} analogs}.
We insert the \gls{mte} analogs to replace the tagging instructions in our tagged heap memory allocator and also modify \thename's instrumentation to automatically insert the \gls{mte} analogs in places of LLVM's \gls{mte} intrinsics.




% \begin{table}[t]
% 	\setlength\tabcolsep{2pt}
% 	\footnotesize
% 	\begin{tabularx}{\columnwidth}{@{}c >{\tt}l Y Y Y >{\hspace{5pt}\tt}l Y Y Y}
% 		\toprule
% 		% \newcommand{\txthead}[1]{\text{\thead{#1}}}
% 		                                                                & \textnormal{\thead{Syscall                                                                                                                                                                                \\ \& API}} & \text{\thead{Base.\\(ns)}} & \text{\thead{Cap.\\(ns)}} & \text{\thead{\underline{Ovh.}\\(\%)}} & \textnormal{\thead{Syscall\\ \& API}} & \text{\thead{Base.\\(ns)}} & \text{\thead{Cap.\\(ns)}} & \text{\thead{\underline{Ovh.}\\(\%)}} \\ \midrule
% 		\multirow{8}{*}[-0.0em]{\thead{\begin{sideways}Syscall\end{sideways}}}
% 		                                                                & socket                                                       & 718.5 & 722.5             & \underline{0.56}     & setsockopt                               & 235.7 & 242.2             & \underline{2.75} \\
% 		% \cmidrule{2-9}
% 		                                                                & \rowcolor[gray]{0.925} bind                                  & 400.0 & 408.0             & \underline{1.98}     & listen                                   & 351.2 & 361.5             & \underline{2.91} \\
% 		% \cmidrule{2-9}
% 		% & accept4 & 1235.9 & 1247.8 & \underline{0.96}
% 		                                                                & accept4                                                      & 1236  & 1248              & \underline{0.96}     & recvfrom                                 & 192.6 & 199.4             & \underline{3.56} \\
% 		% \cmidrule{2-9}
% 		                                                                & \rowcolor[gray]{0.925} openat                                & 471.5 & 486.9             & \underline{3.26}     & read                                     & 247.1 & 254.2             & \underline{2.88} \\
% 		% \cmidrule{2-9}
% 		                                                                & pwrite64                                                     & 320.6 & 328.2             & \underline{2.35}
% 		                                                                & pread64                                                      & 231.8 & 239.0             & \underline{3.06}                                                                                               \\
% 		% \cmidrule{2-9}
% 		                                                                & \rowcolor[gray]{0.925} fstat                                 & 185.9 & 191.5             & \underline{3.03}     & fcntl                                    & 148.1 & 155.8             & \underline{5.20} \\
% 		% \cmidrule{2-9}
% 		                                                                & dup3                                                         & 132.5 & 138.0             & \underline{4.22}     & lseek                                    & 149.7 & 154.8             & \underline{3.38} \\
% 		% \cmidrule{2-9}
% 		                                                                & \rowcolor[gray]{0.925} mmap                                  & 130.0 & 130.0             & \underline{0.01}     & \textbf{getpid}                          & 118.6 & -                 & -                \\ \midrule
% 		\multirow{4}{*}[0em]{\begin{sideways}\thead{API}\end{sideways}} &
% 		\texttt{capac\_malloc}                                          & 23.46                                                        & 50.28 & \underline{114.6} & \texttt{capac\_free} & 14.7                                     & 37.4  & \underline{154.4}                    \\
% 		% \cmidrule{2-9}
% 		                                                                & \rowcolor[gray]{0.925} \texttt{capac\_enter}                 & -     & 167.7             & -                    & \texttt{capac\_exit}                     & -     & 161.2             & -                \\
% 		% \cmidrule{2-9}
% 		                                                                & \texttt{delegate\_ptr}\textsuperscript{*}                    & -     & 327.7             & -                    & \texttt{delegate\_fd}\textsuperscript{*} & -     & 182.9             & -                \\
% 		% \cmidrule{2-9}
% 		                                                                & \rowcolor[gray]{0.925} \texttt{limit\_fd}\textsuperscript{*} & -     & 142.4             & -                    &                                          &       &                   &                  \\
% 		\bottomrule
% 	\end{tabularx}
% 	\begin{tablenotes}
% 		\item \textsuperscript{*} {\small \texttt{capac\_} \ prefix is omitted}
% 	\end{tablenotes}
% 	\caption{The latency of \thename system calls intervened by reference monitor and \libcapac APIs (\textbf{Cap.}), in comparision with the baseline (\textbf{Base.}). Latency of \texttt{getpid} is also measured for comparison.}
% 	\label{tbl:micro}
% 	% \vspace{-2em}
% \end{table}
\begin{table}[t]
	\setlength\tabcolsep{2pt}
	\footnotesize
	\begin{threeparttable}
		\begin{tabularx}{\columnwidth}{@{}c >{\tt}l Y Y Y >{\tt}l Y Y Y}
			\toprule
			 & \thead{Syscall
			\\ \& API} & \thead{Base.\\(ns)} & \thead{Cap.\\(ns)} & \thead{\underline{Ovh.}\\(\%)} &
			\thead{Syscall                                                                                                                                \\ \& API} & \thead{Base.\\(ns)} & \thead{Cap.\\(ns)} & \thead{\underline{Ovh.}\\(\%)} \\
			\midrule
			\multirow{8}{*}[-0.0em]{\thead{\begin{sideways}Syscall\end{sideways}}}
			 & socket                           & 718.5 & 722.5 & \underline{0.56}  & setsockopt                      & 235.7 & 242.2 & \underline{2.75}  \\
			% \rowcolor[gray]{0.925}
			 & bind                             & 400.0 & 408.0 & \underline{1.98}  & listen                          & 351.2 & 361.5 & \underline{2.91}  \\
			 & accept4                          & 1236  & 1248  & \underline{0.96}  & recvfrom                        & 192.6 & 199.4 & \underline{3.56}  \\
			% \rowcolor[gray]{0.925}
			 & openat                           & 471.5 & 486.9 & \underline{3.26}  & read                            & 247.1 & 254.2 & \underline{2.88}  \\
			 & pwrite64                         & 320.6 & 328.2 & \underline{2.35}  & pread64                         & 231.8 & 239.0 & \underline{3.06}  \\
			% \rowcolor[gray]{0.925}
			 & fstat                            & 185.9 & 191.5 & \underline{3.03}  & fcntl                           & 148.1 & 155.8 & \underline{5.20}  \\
			 & dup3                             & 132.5 & 138.0 & \underline{4.22}  & lseek                           & 149.7 & 154.8 & \underline{3.38}  \\
			% \rowcolor[gray]{0.925}
			 & mmap                             & 130.0 & 130.0 & \underline{0.01}  & getpid                          & 118.6 & -     & -                 \\
			\midrule
			\multirow{4}{*}[0em]{\begin{sideways}\thead{API}\end{sideways}}
			 & capac\_malloc                    & 23.46 & 50.28 & \underline{114.6} & capac\_free                     & 14.7  & 37.4  & \underline{154.4} \\
			% \rowcolor[gray]{0.925}
			 & capac\_enter                     & -     & 167.7 & -                 & capac\_exit                     & -     & 161.2 & -                 \\
			 & delegate\_ptr\textsuperscript{*} & -     & 327.7 & -                 & delegate\_fd\textsuperscript{*} & -     & 182.9 & -                 \\
			% \rowcolor[gray]{0.925}
			 & limit\_fd\textsuperscript{*}     & -     & 142.4 & -                 &                                 &       &       &                   \\
			\bottomrule
		\end{tabularx}

		\begin{tablenotes}
			\item \textsuperscript{*} {\small \texttt{capac\_} prefix omitted.}
		\end{tablenotes}

		\caption{Latency of \thename system calls intercepted by the reference monitor and \libcapac APIs (Cap.) vs. baseline (Base.). We also report \texttt{getpid} for context.}
		\label{tbl:micro}
	\end{threeparttable}
\end{table}

% [BENCH] capac_enter 167.713375
% [BENCH] capac_exit 328.89095833333334
% [BENCH] capac_free 14.410083333333333
% [BENCH] delegate_ptr 327.72333333333336
% [BENCH] capac_limit_fd 141.07779166666666
% [BENCH] delegate_fd 182.95304166666668
% [BENCH] capac_alloc 50.28158333333333
% [BENCH] capac_free 37.41125
% [BENCH] dlmalloc 23.462458333333334
% [BENCH] dlfree 14.742291666666667

% 0  &     accept4 &  1235.9 &  1247.8 \\
% 1  &        bind &   400.0 &   408.0 \\
% 2  &        dup3 &   132.5 &   138.0 \\
% 3  &       fcntl &   148.1 &   155.8 \\
% 4  &       fstat &   185.9 &   191.5 \\
% 5  &       ioctl &   165.7 &   697.2 \\
% 6  &      listen &   351.2 &   361.5 \\
% 7  &       lseek &   149.7 &   154.8 \\
% 8  &        mmap &   130.0 &   130.0 \\
% 9  &  newfstatat &   297.2 &   198.7 \\
% 10 &      openat &   471.5 &   486.9 \\
% 11 &     pread64 &   231.8 &   239.0 \\
% 12 &    pwrite64 &   320.6 &   328.2 \\
% 13 &        read &   247.1 &   254.2 \\
% 14 &    recvfrom &   192.6 &   199.4 \\
% 15 &  setsockopt &   235.7 &   242.2 \\
% 16 &      socket &   718.5 &   722.5 \\
% 17 &       write &   337.1 &   345.1 \\



\subsection{Microbenchmarks}
\label{sec:evaluation:mon-microbench}
We perform a set of isolated microbenchmarks on the syscalls intervened by the reference monitor, \libcapac API functions, and the instrumentation. The results are listed in \autoref{tbl:micro}.
% For the intervened syscalls and \thename heap allocator, we compare their latency against their baseline counterparts: syscalls without monitor and unmodified \texttt{dlmalloc}.

\begin{table*}[t]
	\centering
	\footnotesize
	% \newcommand*{\belowrulesepcolor}[1]{%
%   \noalign{%
%     \kern-\belowrulesep
%     \begingroup
%       \color{#1}%
%       \hrule height\belowrulesep
%     \endgroup
%   }%
% }
% \newcommand*{\aboverulesepcolor}[1]{%
%   \noalign{%
%     \begingroup
%       \color{#1}%
%       \hrule height\aboverulesep
%     \endgroup
%     \kern-\aboverulesep
%   }%
% }
\vspace{-8em}
\begin{adjustbox}{angle=90}
\begin{threeparttable}
\scriptsize
    % \begin{tabu} to \textwidth {@{}llp{3cm}ccYcccc@{}}

  \begin{tabularx}{1.6\textwidth}{@{}lp{1.7cm}p{5cm}ccYcccc}
 % \begin{NiceTabular}[width=\textwidth,colortbl-like]{@{}llp{3cm}ccYcccc@{}}
    \toprule
    &\multirow{2}{*}{\thead{Domain}}                          &
    \multirow{2}{*}{\thead{Isolated Objects \{Ref. types\}}} &
    \multirow{2}{*}{\thead{Domain                                                                                                                                                                         \\Switches}} &
    \multirow{2}{*}{\thead{Auth.                                                                                                                                                                          \\Syscalls}} &
    \multirow{2}{*}{\thead{$\mathcal{A}$}} &
    \multicolumn{2}{c}{\thead{Executed PACs}}                &
    \multicolumn{2}{c}{\thead{Executed AUTs}}                                                                                                                                                             \\
    % \cmidrule(lr){7-8} \cmidrule(lr){9-10}
      &   &   &   &   &   & \packey{DB} & \packey{DA} & \autkey{DB} & \autkey{DA} \\[1ex]
    \toprule
    \multirow{4}{*}[-2em]{\begin{sideways}\thead{NGINX + LibreSSL}\end{sideways}}
    &
    % \rule{0pt}{20pt}
    \domconn                                                 & \makecell[l]{Server socket \{\fd\}                                                                                                         \\Client connection socket \{\fd\}}   &
    $2$                                                      & $16$                                                         & $-$ & $-$ & $104$   & $-$         & $590$                                   \rule{0pt}{4.0ex}\\[2ex]
    % \cdashline{2-10}\noalign{\vskip\belowrulesep}
    % \cmidrule{2-10}
    % \aboverulesepcolor{lightgray}
    &  
    \rowcolor[gray]{0.925}
    % \rowcolor{lightgray} \rule{0pt}{20pt}
    \domhs                                                   & \makecell[l]{Server priv. key \{\cpath, \fd, \ptr\}                                                                                        \\TLS session key \{\ptr\}\\Client connection socket \{\fd\textsuperscript{\ddag}\}} &
    $3$                                                      & $18$                                                         & $4$ & $9$ & $10.1$K & $13$        & $134$K                                   \rule{0pt}{6.0ex}\\[2ex]

    
    % \cmidrule{2-10}
    % \cdashline{2-10}\noalign{\vskip\belowrulesep}
    & 
    % \rule{0pt}{20pt}
    
    \domses                                                  & \makecell[l]{TLS session key \{\ptr\textsuperscript{\ddag}\}                                                                               \\Client connection socket \{\fd\textsuperscript{\ddag}\} } &
    \makecell[c]{$7$/$8$/$10$/$13$                                                                                                                                                                        \\/$22$/$37$/$70$\textsuperscript{\dag}} &
    \makecell[c]{$14$/$16$/$20$/$28$                                                                                                                                                                      \\/$44$/$76$/$140$\textsuperscript{\dag}} &
    $3$                                                      &
    \makecell[c]{$43$/$49$/$61$/$85$                                                                                                                                                                      \\/$133$/$229$/$421$\textsuperscript{\dag}} &
    \makecell[c]{$818$/$852$/$917$/$1$K                                                                                                                                                                   \\/$1.3$K/$1.8$K/$2.9$K\textsuperscript{\dag}} &
    \makecell[c]{$3.4$K/$6.5$K/$12$K/$25$K                                                                                                                                                                \\/$50$K/$99$K/$199$K\textsuperscript{\dag}} &
    \makecell[c]{$7.7$K/$10$K/$14$K/$24$K                                                                                                                                                                 \\/$42$K/$79$K/$154$K\textsuperscript{\dag}}   \rule{0pt}{5.0ex}\\[2ex]
    % \cmidrule{2-10}
    % \cdashline{2-10}\noalign{\vskip\belowrulesep}
    & 
    % \rule{0pt}{20pt}
    % \rowcolor{lightgray}
    \rowcolor[gray]{0.925}
    \domamb                                                  & $-$                                                          &
    $-$                                                      &
    \makecell[c]{$94$/$99$/$100$/$109$                                                                                                                                                                    \\/$121$/$145$/$193$\textsuperscript{\dag}}  & $-$ & $-$ & $33.6$K & $-$   & $164$K \rule{0pt}{4.0ex} \\[2ex]
    \midrule
    \multirow{1}{*}[-1.1em]{\begin{sideways}\thead{SSH}\end{sideways}}
    &\textsc{PrivKey}                                                 & \makecell[l]{Client private key \{\cpath,\fd,\ptr\}}   &
    $1$                                                      & $72$                                                         & $55$ & $18$ & $6$K   & $24$         & $33.5$K                                 \rule{0pt}{2.0ex} \\
    % \cmidrule{2-10}
    % \cdashline{2-10}
    &
    % \rowcolor{lightgray}
    \rowcolor[gray]{0.925}
    \domamb                                                  & $-$                                                          &
    $-$                & $37$                                      &
    $-$ & $-$ & $14.7$K & $-$   & $114$K \rule{0pt}{3.0ex}\\[1ex]
    \midrule
    \multirow{2}{*}[-0.9em]{\begin{sideways}\thead{wget}\end{sideways}}
    &\textsc{FileDownload}                                                 & \makecell[l]{Downloaded file \{\cpath,\fd,\ptr\}}   &
    $3$                                                      & $325$                                                         & $1$ & $129$ & $3$K   & $385$         & $8.6$K                                   \rule{0pt}{2.0ex}\\
    % \cmidrule{2-10}
    % \vspace{.1em}
    % \cdashline{2-10}
    % \vspace{.1em}
    &
    \rowcolor[gray]{0.925}
    % \rowcolor{lightgray}
    % \rule{0pt}{20pt}
    \domamb                                                                                                      &
    $-$                                               & $-$ &     $131$ & $-$ & $-$ & $703$ & $-$   & $1.6$K \rule{0pt}{3.0ex}\\[1ex] \bottomrule
  \end{tabularx}
  % \end{tabu}
  % \end{NiceTabular}

  \begin{tablenotes}
    \item  \textsuperscript{\dag}  {\small Measurements from \{16K,32K,64K,128K,256K,512K,1024K\} HTTPS file size configurations } \hspace{1.0em} \textsuperscript{\ddag}  {\small  Delegated references} \vspace{0.5em}
    \item {\small$ \mathcal{A}$ = Number of domain-private memory allocations}
    %\item \textbf{Dom. SW:}\# Domain switches\hspace{0.8em} \textbf{Auth. Syscalls:} \# of \fd-authenticated syscall invocations \hspace{0.8em} \textbf{Dom-Priv MemAlloc.:} \# of \texttt{capac\_malloc()} invocation
    %         \item \hfil \textsuperscript{\dag}  {\small  Delegated references} % \textbf{/} {\small \{Total PAC/AUTs in Program\}}
  \end{tablenotes}

  \caption{\thename domains and protected assets in evaluated applications. Columns 3-9 show runtime measurements of \thename operations executed during a single iteration of \{HTTPS file transfer (NGINX) / SSH handshake / wget file download\}.}

  \label{tbl:porting-app}
  \label{tbl:app}
  % \vspace{-2em}
\end{threeparttable}
\end{adjustbox}
\end{table*}





\hdr{System call latency with reference monitor}
We measured the average latency of \thename-protected syscall used in our evaluated applications that authenticate file paths and \fds, and compared them against un-protected syscalls.
% The evaluated \glspl{syscall} are the ones used by our tested applications.
The overhead induced by \thename includes the latency from syscall hooking and \thename's reference authentication.
Our results reported an average of $2.65\%$ across the measured syscalls.
This shows the efficiency of \thename reference monitor design that fully utilizes hardware acceleration provided by \gls{pa}, not to mention the simplicity of capability where the reference alone is enough to make access control decisions without complex bookkeeping on domains and resources.
Also, this overhead is imposed only on the \thename-enabled processes.




\hdr{\libcapac APIs latencies}
The overhead of \texttt{capac\_malloc()} and \texttt{capac\_free()} originates from domain authentication and \gls{mte}-based memory tagging. 
Measuring the latency of domain private memory allocation and free on blocks of $1$KB, \texttt{capac\_malloc()} and \texttt{capac\_free()} reported approximately $114.6\%$ and $154.4\%$ .
\texttt{capac\_enter} and \texttt{capac\_exit}, \texttt{capac\_delegate\_ptr}, \texttt{capac\_delegate\_fd} and \texttt{capac\_limit\_fd} are simple \texttt{ioctl} calls to the in-kernel reference monitor, where only \texttt{capac\_limit\_fd} does not requires the key switching.
Their latencies are on par with general syscalls.
Besides a roundtrip to the kernel mode, we suspect their main source of overhead is from writing to the \gls{pa} key registers.
% \texttt{capac\_delegate\_ptr} and \texttt{capac\_delegate\_fd} have higher latency due to 
% \texttt{capac\_delegate\_ptr()} and \texttt{delegate\_fd()} whose inner-workings are explained in~\cref{sec:design}, also measured to have reasonable latencies.




\begin{figure}[b]
	\centering
	\includegraphics[width=.95\columnwidth]{images/nbench-overheads.pdf}
	\caption{Overhead of instrumented nbench-byte, normalized to uninstrumented baseline and sorted by overhead.}
	\label{fig:nbench}
\end{figure}





% \FloatBarrier
\hdr{nbench-byte benchmark}
\label{sec:evaluation:nbench}
We measure the isolated overhead of \thename's instrumentation on the \emph{nbench-byte}~\cite{nbench} benchmark, which is also used by several previous works on \gls{pa}-based defenses~\cite{pactight,pac-it-up}.
nbench-byte consists of $11$ CPU and memory subsystem benchmarks that allow us to measure the overhead of \thename in general computations.
\autoref{fig:nbench} shows the results of nbench-byte compiled with \thename's instrumentation that enforces \gls{pa} on on-save signing and on-load authentication of all in-program pointers.
On average, \thename's instrumentation incurs $11.37\%$ performance degradation to the uninstrumented version.
% The highest overhead was observed in FP Emulation (39.20\%) and Huffman (38.48\%).


% \begin{table}
\footnotesize
 \begin{tabularx}{\columnwidth}{@{}lXcl@{}}
 \toprule
 \thead{Applications} & \thead{Isolated Objects \{Ref. Types\}}& \thead{No. of\\domains}& \thead{Modified\\SLOC} \\
 \midrule
 \makecell[l]{NGINX (v1.23.0)\\+LibreSSL (vx,x,x)} & \makecell[l]{Private key\{\cpath, \fd,\ptr\}\\TLS Session key\{\ptr\}\\Sokets \{\fd\}} & 3 \\
 \midrule
 OpenSSH (v1.1.1) & Private key \{\cpath, \fd,\ptr\} & 1 \\
 \midrule
 minizip &Password \{\ptr\}& 1\\
 \midrule
 wget (v1.1.1) & Password \{\ptr\}& 1\\

 
 \bottomrule
 \end{tabularx}

\caption{Programs used to evaluate \thename's resource isolation. }

\label{tbl:porting-app}
\label{tbl:app}
\end{table}




%We expect the reason is that our instrumentation also works with compiler optimizations, which removed many redundant on-load authentications, while PARTS disables all optimizations in their tests~\cite{pac-it-up}. --> discuss this separately maybe design:instrumentaiton.




% \vspace{-1em}
\subsection{Adapting \thename to real-world programs}
We exemplify \thename's capability domains by adapting it to several open-source applications studied by previous isolation frameworks~\cite{erim,shreds,ptrsplit}. 
These applications include OpenSSH's \texttt{ssh} utility~\cite{openssh}, \texttt{wget} file download utility~\cite{wget}, and NGINX webserver~\cite{nginx} compiled along the TLS library LibreSSL~\cite{libressl}. 
Representative code examples of our modifications can be found in the appendix of the extended version of this paper~\cite{capacity-extended}.
% Different from other compartmentalization frameworks~\cite{erim,jenny,shreds,ptrsplit,hakc}, on top of memory isolation, \thename also enables the coherent in-process isolation of system resources, such as file descriptors of private files and sockets.
The applications are also compiled with the \thename-instrumented \emph{musl} libc~\cite{musl} adapted to handle signed pointers in system call arguments.
\autoref{tbl:porting-app} summarizes the \thename domains in protected applications and their protected assets.
%Our instrumentation framework is proven mature enough to compile complex real-world programs such as the webserver and its components.
% We report the performance results for ssh in this section.


% More specifically, we measured the cumulative occurrences of domain switches, \fd-authenticated \gls{syscall} invocation, \texttt{capac\_malloc()} invocation, and \texttt{pac}/\texttt{aut} instructions \emph{executed} for the varying file sizes.




\hdr{OpenSSH's ssh (v9.3p1)}
% The \texttt{ssh} utility uses a private key from the file system to establish a secure session with the remote host. 
We demonstrate \thename's program secrets protection throughout their life-cycles with \texttt{ssh}.    
We compile OpenSSH to use its built-in crypto library and use \thename to isolate the private key of \texttt{ssh}.
We introduce a new domain, \textsc{Privkey}, with exclusive access to the private key file and its in-memory buffer.
This domain spans from when the private key is loaded from the file system until after the login operation succeeds.
 % Unlike previous systems \cite{shreds,ptrsplit}, \thename also protects the path and file descriptor of the private key.
% To achieve this, 
We then assign the static private key to the domain, allocate the private key buffer (\texttt{struct sshkey}) in private memory, annotate its pointers with \texttt{DOM\_PRIV}, and finally wrap functions that access the private key within domain entry gates.
These changes require about $50$ LoC changes over 13 functions across 5 files.
% in the files \texttt{sshkey.c}, \texttt{authfile.c}, \texttt{sshbuf.c}, \texttt{sshconnect2} and \texttt{ssh.c}. 


\begin{figure*}[t]
	\centering
	% \includegraphics[width=0.92\textwidth]{images/webserver-capacity-final.pdf}
	\includegraphics[width=\textwidth]{images/webserver-overheads-3.pdf}
	\begin{minipage}[t]{.49\linewidth}
		\centering
		% \vspace{-1em}
		\subcaption{Single-threaded}
		 \vspace{-1em}
	\end{minipage}%
	\begin{minipage}[t]{.49\linewidth}
		\centering
		% \vspace{-1em}
		\subcaption{Multi-threaded (8 Threads)}
		 \vspace{-1em}
	\end{minipage}
	\caption{\textsc{Baseline} webserver vs. \thename-enabled webserver throughput benchmark performed with local \texttt{ab} client.
		% The \texttt{Baseline} is unmodified NGINX, whereas the \thename-enabled webserver runs with \thename reference monitor, \libcapac, and instrumented software components. 
	}
	\label{tab:exec-macro}
\end{figure*}


\hdr{\texttt{wget}  (v1.21.2)}
% We port \texttt{wget} 
We use \thename to isolate the file received from the internet in \texttt{wget} to demonstrate least-privilege compartmentalization.
We create a domain called \textsc{FileDownload} with exclusive access to the downloaded buffer and the output file path/\fd. 
We then modify the initialization of \texttt{wget} to assign the output file access to \textsc{FileDownload}, and use \texttt{cap\_limit\_fd} to revoke all other attributes from the file's \fd except \texttt{CAP\_WRITE}. 
This ensures that the domain has write-only access to the output \fd \footnote{While \texttt{wget} already make the file handle write-only, \thename enables more fine-grained in-process access control that supports read and write domains.}.
The domain encapsulates the function \texttt{retrieve\_url}, which fetches the requested file from a URL into a \texttt{capac\_malloc}-allocated buffer before writing into the output file's \fd. 
We also annotate pointers variables that point to the downloaded file in memory with \texttt{CAPAC\_VAR}. 
The total modifications to \texttt{wget} is about $30$ LoC.



\begin{figure}[b]
	% \centering	\includegraphics[width=0.9\columnwidth]{images/webserver-domains-libertine.pdf}
	\centering	\includegraphics[width=0.98\columnwidth]{images/nginx-domains.pdf}
	\caption{\thename enables life-cycle private key protection and per-instance isolation of session keys in the NGINX webserver.}
	\label{fig:webserver-domains}
\end{figure}



\label{sec:using-capacity}\label{sec:design:webserver}\label{sec:evaluation:porting-app}\label{sec:usage}




\hdr{NGINX (v1.23.0) + LibreSSL (v3.5.3)}
We use the NGINX webserver and its crypto library, LibreSSL, to demonstrate the multi-domain interaction of \thename. 
We create three domains in the webserver, namely \domhs, \domses, and \domconn to isolate the server private key, session keys, and connection sockets with the least privilege principle in mind.
\autoref{tbl:porting-app} displays a summary of the webserver domains and their isolated resources.
We modified $223$ LoC in LibreSSL and $114$ LoC in NGINX.



% The rest of the webserver not in a domain is the ambient domain, \domamb.
% We construct the above domains by using the approach described in \cref{sec:design:prog-model}.
% Then, the webserver and its dependencies are compiled with \thename instrumentation framework.
% for complete mediation of pointer uses as well as domain-aware \gls{pa} instruction insertion.



\autoref{fig:webserver-domains} illustrates the main operations and their interaction with domain-private objects in domain \domhs and \domses.
\domhs, which consists of $5$ NGINX and $2$ LibreSSL functions endowed with exclusive rights to the server private key.


% \deleted{
	% In \thename's design, this means that all object references that can point to the private key in its life-cycle are signed and authenticated with \domhs's domain key instead of the ambient key.
% }
% Hence, the context must be inside \domhs to pass the authentication and interact with private key objects.
% The private key protection is bootstrapped by \texttt{capac\_init()} that requests the kernel to bind the private key's file-system path to its owner domain during process initialization.
% \deleted{
% 	From then on, the reference monitor would only allow a thread to open (\texttt{sys\_open}) or read (\texttt{sys\_read}) if and only if the thread is executing inside \domhs and \fds carry a valid PAC with the domain's key.
% }
% \st{In \domhs, }
% Additionally, the in-memory objects that hold the contents from the private key file are allocated with \texttt{capac\_malloc()}.
% Additionally, the private key-containing structure (\texttt{struct EVP\_PKEY}) is placed in the domain-tagged memory with \texttt{capac\_malloc()}/, and its pointers are applied with \texttt{DOM\_PRIV}.
% \texttt{DOM\_PRIV} is also applied to the \texttt{struct* EVP\_PKEY} arguments of functions that interact with the now-in-memory private key.
\domses protects session keys by hosting $17$ functions,  mostly located inside LibreSSL's AES-GCM implementation.
We generate a unique instance ID and use it as the modifier for each session since a session's resources are mostly temporary and can be isolated. 
% Since the session key is generated anew on each connection through SSL handshaking, a connection-generated modifier isolates the SSL key's pointer, ensuring can only decrypt messages for a specific connection.
As \domhs execution finishes, it produces a TLS session key, which is to be \emph{delegated} to the corresponding \domses instance using the API \texttt{capac\_delegate\_ptr} (\autoref{tab:apis}).
% After delegation, the session key is in \domses's domain-private memory, and the domain is received a domain-privated pointer to the object (\texttt{struct* aead\_aes\_gcm\_ctx}).
The domain-private stack  is a noteworthy feature of \thename used in \domses.
When a function performs complex cryptographic operations, we use \texttt{DOM\_PRIV\_STACK} instead of tracking the data flow of sensitive data bouncing around in the function's local variables.
% \hdr{Socket delegation and isolation} 



Finally, \domconn manages the webserver's sockets.
We wrap the functions that handle sockets inside its domain switches of \domconn, such that when the server requests its \emph{server socket} \fd using \texttt{sys\_socket}, the returned \fd is domain-private.
On client connection requests, \domconn invokes \texttt{sys\_accept4} with the server socket as the input and retrieves a \emph{client} socket \fd.
\domconn then delegate the returned \emph{client} socket \fd to \domhs, using \texttt{capac\_delegate\_fd}. 
After TLS handshaking, \domhs again delegates the socket \fd to \domses.
Notably, \texttt{CAP\_SOCKET} attribute is enabled by default when socket \fds are created, segregating socket and non-socket \fds in syscalls without the user's intervention.



% \vspace{-.5}


% \begin{tabular}{cc}
% \toprule
%        \thead{Program} &\thead{Overhead(\%)}  \\
%         \midrule
%      \texttt{ssh} & $0.37$ \\
%      \texttt{wget} & $5.95$ \\
%      \bottomrule
% \end{tabular}

% \begin{floatingfigure}{0.5\columnwidth}
% \footnotesize
% \centering

% \caption{test}
% \end{floatingfigure}


\begin{table}[t]
    \centering
% \footnotesize
\small
\begin{tabular}{lccc}
\toprule
       \thead{Program} &\thead{Baseline (ms)} &\thead{w/ \thename (ms)}&\thead{Overhead (\%)}  \\
        \midrule
     \texttt{ssh}&$347.35$ & $348.62$  & $0.37$ \\
 \rowcolor[gray]{0.925}    \texttt{wget}&  $17.12$ & $18.14$&$5.95$ \\
     \bottomrule
\end{tabular}
    \caption{\thename overheads on  \texttt{ssh} and \texttt{wget}.}
    \label{tbl:sshwget}
\end{table}


\subsection{Application performance benchmark} \label{sec:evaluation:macrobench}
We perform performance benchmarks on the adapted applications.
These applications run with all design components of \thename explained thus far, including compile-time instrumentation, compatibility fixes, \libcapac, the kernel-level reference monitor, and backward and forward-edge CFI.
To accurately profile the source of \thename's performance overhead, we inserted probes into the reference monitor, \libcapac, and the instrumented program code to collect \emph{runtime} execution count of \thename operations during a single iteration of execution, as shown in~\autoref{tbl:porting-app}.


% \begingroup
% \setlength{\intextsep}{5pt}%
% \setlength{\columnsep}{10pt}%
% \begin{wrapfigure}{r}{3.4cm}

% \caption{\thename overhead on \texttt{ssh} and \texttt{wget}.}
% \label{tbl:sshwget}
% % \vspace{-3mm}
% \end{wrapfigure}
% \hdr{Performance overheads}
\hdr{\texttt{ssh} and \texttt{wget}} 
For \texttt{ssh}, we measure the average overhead of connecting to a local server $10,000$ times and compare it against the unmodified version.
We evaluate the protected \texttt{wget} by timing the average overhead in latency when requesting a 1MB file from a local http webserver $10,000$ times compared to the baseline.
\autoref{tbl:sshwget} displays the \thename-incured performance overheads in \texttt{ssh} and \texttt{wget}, which shows only minimal slowdown for both applications. 


% \endgroup

\hdr{NGINX + LibreSSL} We evaluate the throughput of the \thename-enabled NGINX webserver against its unmodified counterpart, which we shall call \textsc{Baseline}.
% The isolated domains and their assets are as previously explained in \cref{sec:usage}.
We use the Apache benchmarking tool, \texttt{ab}~\cite{apache-benchmark} locally (i.e., no network latency) to measure the number of requests processed per second, with varying file sizes, \{$16$K, $32$K ... $1$M\}, on also single-threaded and multi-threaded (8 threads) settings.
$100,000$ requests were performed, and the keep-alive option was used to avoid re-initializing the connections.
The cipher suite negotiated between the client and the server was \texttt{ECDHE-RSA-AES256-GCM-SHA384}.


The webserver benchmark results are displayed in \autoref{tab:exec-macro}, which shows the average performance degrades about $17\%$ for the single-threaded server and $13.54\%$ for the multi-threaded settings across all file size configurations.
The latency measurements for \thename operations from our microbenchmark (\autoref{tbl:micro}) allow us to reason that the performance overheads from the reference monitor and \libcapac would be very limited.
The varied file sizes allow us to capture the characteristics of \thename's performance overhead, which increases as the requested file sizes increase.
Along with the runtime statistic collected in \autoref{tbl:app}, it can be observed that \thename's overhead protecting pointers and memory is proportional to the amount of computation.
We can conclude that the performance overhead of \thename would be predominantly from the increased number of instructions due to instrumentation.














% Note that API functions that are unique to \thename (e.g., \texttt{capac\_enter()}) do not have counterparts.
% The reference monitor interposes a set of designated syscalls and enforces an efficient and minimal authentication routine centered around the \gls{pa} instructions, as described in~\cref{sec:implementation:ref-monitor}.








%\hdr{Latency of \texttt{pac} and \texttt{aut} instructions}
%We also measure the execution time of individual \texttt{pac} and \texttt{aut} instructions. We found that all \texttt{pac} and \texttt{aut}-family instruction performs roughly the same, about $13.1$ns per instruction.
